% chapters/chapter3.tex
\section{РЕАЛІЗАЦІЯ ПРОГРАМНОГО РІШЕННЯ}

% Розділ 3 згідно нових методичних вказівок (15-20 стор.):
% Для типу роботи «ПЗ / Веб / ІС» підрозділи 3.4 «Дані та їх підготовка»
% (обов'язковий лише для ML/ШІ) та 3.5 «Апаратна частина» (обов'язковий
% лише для IoT / embedded) не застосовуються. Використано послідовну
% нумерацію без пропусків:
%
% 3.1 Опис технологій та інструментів розробки
% 3.2 Структура проєкту та організація вихідного коду
% 3.3 Реалізація ключових модулів та алгоритмів
% 3.4 Обробка виняткових ситуацій та забезпечення надійності
% 3.5 Розгортання системи
% 3.6 Тестування програмного забезпечення
% 3.7 Висновки до розділу 3

% -------------------------------------------------
\subsection{Опис технологій та інструментів розробки}

У підрозд.~2.1 обґрунтовано архітектурні рішення: клієнтський SPA,
реактивний рендеринг, централізоване сховище стану, спеціалізована
бібліотека графової візуалізації. У цьому підрозділі наведено
конкретні технології, бібліотеки, їхні версії та середовище розробки,
що реалізують зазначені архітектурні рішення.

\subsubsection*{Мови програмування та виконуване середовище}

Основною мовою реалізації обрано \textbf{TypeScript~5.x} --- надбудову
над JavaScript зі статичною типізацією. Мова компілюється у
стандартизований \textbf{ECMAScript~2020 (ES11)}, який виконується
усіма сучасними браузерами без додаткових polyfill'ів, що
задовольняє вимогу NF4. Статична типізація дозволяє виявляти на
етапі компіляції помилки передачі об'єкта \texttt{Token} замість
числа чи неправильної форми комплексного значення, що критично для
дотримання NF1 (точність $\varepsilon \le 10^{-9}$).

Цільовим середовищем виконання є браузерний рушій JavaScript
(V8 у Chrome/Edge, SpiderMonkey у Firefox) з підтримкою
\texttt{requestAnimationFrame}, Web Workers, Canvas API та Local
Storage.

\subsubsection*{Фреймворки та бібліотеки}

Повний перелік використаних бібліотек із версіями та призначенням
наведений у табл.~\ref{tab:tech_stack}.

\begin{table}[H]
  \caption{Технологічний стек реалізації}
  \label{tab:tech_stack}
  \centering
  \small
  \begin{tabular}{p{3.2cm}p{1.6cm}p{8.2cm}}
    \toprule
    Технологія & Версія & Призначення \\
    \midrule
    TypeScript & 5.x & Статично типізована мова реалізації \\
    Next.js & 14.x & React-фреймворк, статичний експорт SPA \\
    React & 18.x & Реактивний рендеринг UI-компонентів \\
    React Flow (\texttt{@xyflow/react}) & 12.x & Інтерактивна графова візуалізація DFG \\
    Zustand & 4.x & Централізоване реактивне сховище стану \\
    Comlink & 4.x & RPC-обгортка для Web Workers \\
    fft.js & 4.x & Референтна реалізація ДПФ для верифікації \\
    Jest & 29.x & Фреймворк модульного тестування \\
    ESLint & 8.x & Статичний аналіз коду, правила стилю \\
    Prettier & 3.x & Автоматичне форматування коду \\
    \bottomrule
  \end{tabular}
\end{table}

\textbf{Next.js} забезпечує маршрутизацію, оптимізоване завантаження
ресурсів і статичну збірку без серверної частини (режим
\texttt{output: 'export'}). Усі сторінки симулятора позначені
директивою \texttt{'use client'}, оскільки синхронне клієнтське
виконання \texttt{requestAnimationFrame}-циклу несумісне зі
server-side rendering.

\textbf{React Flow} надає готові механізми перетягування вузлів,
масштабування, анімованих ребер і кастомних типів вузлів. Його
абстрактна модель \texttt{Node}/\texttt{Edge} безпосередньо
відображається на формальну модель DFG (підрозд.~2.4).

\textbf{Zustand} --- легка альтернатива Redux без шаблонного коду
редьюсерів. Забезпечує реактивне сховище стану, до якого компоненти
підписуються через хук \texttt{useSimStore(selector)} з мемоїзацією.

\textbf{fft.js} --- компактна JavaScript-реалізація ДПФ, що
використовується виключно як референтна для верифікації
(підрозд.~3.6) і не залучається до шляху виконання основного
алгоритму, що забезпечує незалежність результату симулятора від
зовнішньої бібліотеки.

\subsubsection*{Середовище розробки та інструментарій}

Розробка виконується у редакторі \textbf{Visual Studio Code}
з розширеннями ESLint, Prettier і TypeScript. Pre-commit-хуки
реалізовано через \texttt{husky} + \texttt{lint-staged}: перед
кожним коммітом автоматично запускаються ESLint та Prettier
для змінених файлів.

Менеджер пакетів --- \textbf{pnpm}, що забезпечує швидке встановлення
залежностей і ефективне використання дискового простору через
hard-linking. Файлом залежностей є \texttt{package.json}; файл
\texttt{pnpm-lock.yaml} гарантує детермінованість збірки на різних
машинах.

\subsubsection*{Система керування версіями та репозиторій}

Використовується \textbf{Git} зі стратегією \textit{trunk-based
development}: основна розробка ведеться у гілці \texttt{main}, для
великих змін створюються короткоживучі feature-гілки. Репозиторій
розміщений на GitHub, що забезпечує доступність для науково-педагогічного
комітету БКР.

Опційно налаштовано \textbf{GitHub Actions} для автоматичного запуску
модульних тестів і лінтерів на кожному push до \texttt{main}
(див. підрозд.~3.5).


% -------------------------------------------------
\subsection{Структура проєкту та організація вихідного коду}

\subsubsection*{Верхньорівневе дерево репозиторію}

Структура кореневого каталогу проєкту організована за стандартом
Next.js з додатковими директоріями для тестів і документації
(рис.~\ref{fig:project_tree}).

\begin{figure}[H]
\begin{verbatim}
simulator/
+-- public/                  # статичні ресурси (SVG-іконки, favicons)
+-- src/
|   +-- app/                 # Next.js App Router
|   |   +-- layout.tsx
|   |   \-- page.tsx
|   +-- core/                # ядро симуляції (незалежне від UI)
|   |   +-- types.ts
|   |   +-- scheduler.ts
|   |   \-- buffers.ts
|   +-- graph/               # генерація та візуалізація графа
|   |   +-- generateDFT4x4.ts
|   |   +-- layout4x4.ts
|   |   +-- GraphView.tsx
|   |   \-- TokenEdge.tsx
|   +-- components/          # UI-компоненти
|   |   +-- SpectrumView.tsx
|   |   +-- Metrics.tsx
|   |   +-- Controls.tsx
|   |   +-- TerminalOutput.tsx
|   |   \-- nodes/
|   |       +-- SourceNode.tsx
|   |       +-- SinkNode.tsx
|   |       +-- DFT4Node.tsx
|   |       \-- TwiddleNode.tsx
|   +-- store/               # централізоване сховище стану
|   |   \-- simStore.ts
|   +-- utils/               # допоміжні утиліти
|   |   +-- math.ts
|   |   \-- compare.ts
|   \-- workers/             # Web Workers
|       \-- fft.worker.ts
+-- tests/                   # Jest-тести
|   +-- math.test.ts
|   +-- scheduler.test.ts
|   \-- compare.test.ts
+-- package.json
+-- tsconfig.json
+-- next.config.js
\-- .eslintrc.js
\end{verbatim}
\caption{Структура директорій проєкту. Джерело: розроблено автором.}
\label{fig:project_tree}
\end{figure}

\subsubsection*{Принцип декомпозиції}

Декомпозиція на модулі виконана за принципом \emph{розподілу за
шарами відповідальності}, що відповідає архітектурі Container з
підрозд.~2.3:

\begin{itemize}
  \item \texttt{src/core/} --- ядро симуляції без будь-яких
        залежностей від UI; може бути імпортоване у Node.js для
        batch-обчислень і використовується в unit-тестах;
  \item \texttt{src/graph/} --- генерація структури DFG і компоненти
        інтерактивної візуалізації на React Flow;
  \item \texttt{src/components/} --- UI-компоненти, що не залежать
        від React Flow (панелі керування, метрики, спектр);
  \item \texttt{src/store/} --- реактивне сховище стану та
        інтеграція з ядром симуляції;
  \item \texttt{src/utils/} --- чисті функції (математика, верифікація)
        без побічних ефектів;
  \item \texttt{src/workers/} --- Web Worker для фонових обчислень.
\end{itemize}

Така організація забезпечує можливість ізольованого тестування ядра
симуляції без браузерного середовища та спрощує пошук місць
внесення змін при підтримці коду.

\subsubsection*{Відповідність компонентів архітектури та модулів коду}

Таблиця~\ref{tab:arch_to_code} формалізує зв'язок між компонентами
архітектури з підрозд.~2.3 та модулями вихідного коду.

\begin{table}[H]
  \caption{Відповідність «Компонент $\to$ Модуль/файл»}
  \label{tab:arch_to_code}
  \centering
  \small
  \begin{tabular}{p{3.5cm}p{5.5cm}p{4.2cm}}
    \toprule
    Компонент архітектури (2.3) & Модуль / файл & Роль \\
    \midrule
    Simulation Core & \texttt{src/core/scheduler.ts} & Клас \texttt{DataflowRuntime}, метод \texttt{tick()} \\
    Simulation Core & \texttt{src/core/buffers.ts} & FIFO-черги на дугах \\
    Simulation Core & \texttt{src/core/types.ts} & Типи \texttt{Token}, \texttt{Edge}, \texttt{NodeSpec} \\
    Graph Generator & \texttt{src/graph/generateDFT4x4.ts} & Побудова DFG \\
    Graph Generator & \texttt{src/graph/layout4x4.ts} & Координати вузлів \\
    State Store & \texttt{src/store/simStore.ts} & Zustand-сховище \\
    UI Layer (GraphView) & \texttt{src/graph/GraphView.tsx} & React Flow + rAF-цикл \\
    UI Layer (SpectrumView) & \texttt{src/components/SpectrumView.tsx} & Canvas-діаграма $|X(k)|$ \\
    UI Layer (Metrics) & \texttt{src/components/Metrics.tsx} & EMA-метрики \\
    UI Layer (Controls) & \texttt{src/components/Controls.tsx} & Пуск/пауза/швидкість \\
    UI Layer (вузли) & \texttt{src/components/nodes/*.tsx} & SourceNode, SinkNode, DFT4Node, TwiddleNode \\
    Math \& verify & \texttt{src/utils/math.ts}, \texttt{compare.ts} & Комплексна арифметика, верифікація \\
    \bottomrule
  \end{tabular}
\end{table}

\subsubsection*{Ключові конфігураційні файли}

\texttt{tsconfig.json} визначає налаштування компілятора TypeScript:
увімкнені \texttt{strict}, \texttt{noImplicitAny},
\texttt{strictNullChecks} --- для максимальної безпеки типів.
Цільовий стандарт --- ES2020.

\texttt{next.config.js} встановлює \texttt{output: 'export'} для
статичної збірки та налаштовує базовий шлях. Функцій і API немає,
усі сторінки клієнтські.

\texttt{.eslintrc.js} налаштовує правила ESLint з розширеннями
\texttt{@typescript-eslint}, \texttt{react-hooks},
\texttt{@next/next}. Заборонено \texttt{any}, \texttt{console.log}
у production, не імпортовані змінні, побічні ефекти у хуках
без залежностей.

\subsubsection*{Стандарти кодування}

Код документується у стилі \textbf{JSDoc}, з особливим акцентом
на публічних методах \texttt{DataflowRuntime}, \texttt{simStore} і
експортованих утилітах. Типи та інтерфейси описуються безпосередньо
у \texttt{types.ts}. Ідентифікатори --- англомовні, камел-кейс для
змінних і функцій, паскаль-кейс для типів і React-компонентів.
Максимальна довжина рядка --- 100~символів, ця межа автоматично
контролюється Prettier.


% -------------------------------------------------
\subsection{Реалізація ключових модулів та алгоритмів}

\subsubsection*{Клас DataflowRuntime}

Клас \texttt{DataflowRuntime} (модуль \texttt{src/core/\allowbreak scheduler.ts})
є програмною реалізацією абстрактної машини потоку даних
з підрозд.~2.2~\cite{veen1986,dennis1975}. Він відповідає за
виконання правила активації вузлів та управління FIFO-чергами
на дугах графа.

Клас зберігає такі внутрішні структури даних:
\begin{itemize}
  \item незмінне посилання на граф \texttt{Graph} (вузли і дуги);
  \item відображення \texttt{Map<edgeId, FIFO<Token>>} --- окрема
        черга токенів для кожної дуги;
  \item \texttt{readyAt: Map<nodeId, number>} --- час, раніше якого
        вузол не може активуватись повторно (моделює \texttt{latency});
  \item \texttt{inbox: Map<nodeId, Map<portId, Token>>} --- токени,
        що вже доставлені до вузла, але ще не поглинуті;
  \item обробники подій: \texttt{onToken}, \texttt{onFire},
        \texttt{onBeforeFire}, \texttt{onOutput}.
\end{itemize}

Логічний час надається ззовні через функцію
\texttt{getTime: () => number}, що передається у конструктор. У
симуляторі цією функцією є \texttt{() => store.simTime}, завдяки чому
годинник ядра повністю відокремлений від астрономічного часу.

\subsubsection*{Метод tick() --- алгоритм одного кроку симуляції}

Метод \texttt{tick(dt, maxFires)} виконує єдиний крок потокового
обчислення і реалізує правило активації (firing rule), описане у
підрозд.~2.2~\cite{veen1986,dennis1975}. Алгоритм складається
з чотирьох послідовних кроків:

\begin{enumerate}
  \item \textbf{Доставка токенів.} Для кожної дуги $e$ перевіряється,
    чи найстарший токен у черзі $\mathrm{buf}[e.\mathrm{id}]$
    задовольняє умові доставки:
    \begin{equation}
      t_{\text{now}} - \mathrm{token}.t \;\geq\; e.\mathrm{delay},
      \label{eq:delivery}
    \end{equation}
    де $t_{\text{now}} = \texttt{simTime}$. Якщо умова~(\ref{eq:delivery})
    виконана, токен видаляється з черги і поміщається у \texttt{inbox}
    відповідного вузла-приймача; емітується подія \texttt{onToken}.

  \item \textbf{Перевірка готовності вузла.} Вузол $n$ є
    \emph{готовим} до активації, якщо одночасно виконуються дві умови:
    \begin{itemize}
      \item[(а)] у \texttt{inbox}$[n]$ є токен на \emph{кожному}
        вхідному порті: $\forall p \in n.\mathrm{inPorts}:
        \texttt{inbox}[n].\mathrm{has}(p)$;
      \item[(б)] вузол вільний від блокування:
        $\texttt{readyAt}[n] \leq t_{\text{now}}$.
    \end{itemize}

  \item \textbf{Активація вузла.} Для кожного готового вузла
    (не більше \texttt{maxFires} за один виклик):
    \begin{itemize}
      \item[3.1] токени поглинаються з \texttt{inbox}$[n]$ і
        передаються у метод \texttt{execute(node, inputs)};
      \item[3.2] виконується логіка вузла за типом
        \texttt{NodeKind}, обчислюються вихідні токени;
      \item[3.3] час блокування оновлюється:
        $\texttt{readyAt}[n] \leftarrow t_{\text{now}} +
        n.\mathrm{latency}$;
      \item[3.4] емітуються події \texttt{onBeforeFire} і
        \texttt{onFire}.
    \end{itemize}

  \item \textbf{Розміщення вихідних токенів.} Для кожного вихідного
    токена знаходяться всі дуги з відповідним портом-джерелом; токен
    поміщається у чергу кожної такої дуги з часом народження:
    \begin{equation}
      \mathrm{token}.t = t_{\text{now}} + n.\mathrm{latency},
      \label{eq:emit_time}
    \end{equation}
    що гарантує правильну затримку доставки на наступному кроці.
\end{enumerate}

Такий порядок кроків забезпечує детермінованість: при фіксованому
порядку обходу вузлів і дуг результат симуляції не залежить від
реального часу виконання у браузері (NF5).

\subsubsection*{FIFO-черга (buffers.ts)}

Клас \texttt{FIFO<T>} (модуль \texttt{src/core/buffers.ts}) ---
узагальнена черга на базі масиву з методами \texttt{push(item)},
\texttt{pop(): T | undefined}, \texttt{peek(): T | undefined} та
властивістю \texttt{size: number}. Одна черга ставиться у
відповідність кожній дузі графа; черги ініціалізуються у конструкторі
\texttt{DataflowRuntime} ітерацією по \texttt{graph.edges}.

Асимптотична складність операцій: \texttt{push} --- $O(1)$
амортизований, \texttt{pop} --- $O(n)$ у найгіршому випадку (через
зсув масиву). При характерних для симулятора розмірах черг
($n \le 16$) це не становить проблеми; для масштабування на більші
$N$ заплановано перехід на циклічний буфер з $O(1)$-операціями.

\subsubsection*{Логічний час і speed multiplier}

Логічний час \texttt{simTime} збільшується методом
\texttt{tickSimTime()} пропорційно до астрономічного часу та
параметра \texttt{speed}:

\begin{equation}
  \texttt{simTime}_{i+1} = \texttt{simTime}_{i} +
    \Delta t_{\text{wall}} \cdot \texttt{speed},
  \label{eq:simtime}
\end{equation}

де $\Delta t_{\text{wall}}$ --- реальний проміжок між кадрами
(у мілісекундах), а \texttt{speed} --- множник у діапазоні
$[0{,}1;\; 10]$. Відокремлення логічного часу від астрономічного
дозволяє досліджувати поведінку алгоритму у довільному
масштабі без зміни логіки активації.

\subsubsection*{Реалізація вузлів за типом}

Метод \texttt{execute(node, inputs)} реалізує шаблон Strategy
(підрозд.~2.1) через перемикач за \texttt{node.kind}.

\textbf{SOURCE.} Вузол-джерело не має вхідних портів; його токени
вводяться ззовні через виклик
\texttt{runtime.emitFrom(nodeId, 'out', token)}. Токен містить
значення $x(n) = \{\,re, im\,\}$ і поля
$\texttt{t} = \texttt{originT} = \texttt{simTime}$.

\textbf{DFT4.} Вузол приймає 4~токени через порти \texttt{in0}
\ldots \texttt{in3} і застосовує алгоритм метелика
radix-4~(\ref{eq:dft4_butterfly}) з підрозд.~2.2. Лістинг~\ref{lst:dft4}
наводить реалізацію:

\begin{figure}[H]
\begin{verbatim}
function computeDFT4(x: Complex[4]): Complex[4] {
  const E0 = cplxAdd(x[0], x[2]);   // x0 + x2
  const E1 = cplxSub(x[0], x[2]);   // x0 - x2
  const O0 = cplxAdd(x[1], x[3]);   // x1 + x3
  const O1 = cplxSub(x[1], x[3]);   // x1 - x3
  return [
    cplxAdd(E0, O0),                // X0 = E0 + O0
    cplxSub(E1, cplxMulJ(O1)),      // X1 = E1 - j*O1
    cplxSub(E0, O0),                // X2 = E0 - O0
    cplxAdd(E1, cplxMulJ(O1))       // X3 = E1 + j*O1
  ];
}
\end{verbatim}
\caption{Реалізація 4-точкового ДПФ-метелика без нетривіальних
         множень. Функція \texttt{cplxMulJ(z) = \{-z.im, z.re\}}
         реалізує множення на $j$ через перестановку. Джерело:
         розроблено автором.}
\label{lst:dft4}
\end{figure}

Вузол породжує 4~вихідних токени через порти
\texttt{out0} \ldots \texttt{out3}. Жодного нетривіального множення
не виконується.

\textbf{TWIDDLE.} Вузол приймає 1~токен і виконує комплексне
множення на поворотний множник $W_{16}^k = e^{-j 2\pi k/16}$, де
ціле число $k$ зберігається у \texttt{NodeSpec.params.twiddle.k}.
Формули обчислення:

\begin{equation}
  \begin{aligned}
    \mathrm{re}' &= z.\mathrm{re} \cos\theta - z.\mathrm{im} \sin\theta, \\
    \mathrm{im}' &= z.\mathrm{re} \sin\theta + z.\mathrm{im} \cos\theta,
  \end{aligned}
  \quad \theta = -\frac{2\pi k}{16}.
  \label{eq:twiddle_impl}
\end{equation}

При $k = 0$ маємо $W_{16}^0 = 1$, і множення є тривіальним ---
вхідний токен передається на вихід без змін.

\textbf{SINK.} Вузол приймає 1~токен, передає його значення у
\texttt{store.sinks[nodeId]} через серіалізацію
\texttt{JSON.stringify} і емітує подію \texttt{onOutput}.
Одночасно обчислюється наскрізна затримка:

\begin{equation}
  \mathrm{latency} = \texttt{simTime} - \mathrm{token}.\texttt{originT},
  \label{eq:lat_calc}
\end{equation}

і оновлюється EMA-метрика \texttt{latencyEmaMs}.

\subsubsection*{Функція generateDFT4x4()}

Функція \texttt{generateDFT4x4(latency): Graph} у модулі
\texttt{src/graph/\allowbreak generateDFT4x4.ts} реалізує шаблон Factory
(підрозд.~2.1) і будує DFG за алгоритмом:

\begin{enumerate}
  \item створюються 16~вузлів \texttt{source}
        (\texttt{src0} \ldots \texttt{src15}), кожен з одним
        вихідним портом \texttt{out};
  \item створюються 4~вузли \texttt{dft4} першого ярусу
        (\texttt{row-0} \ldots \texttt{row-3}) з 4~входами та
        4~виходами і \texttt{latency} $= 400$;
  \item додаються 16~дуг \texttt{src}$_i \to$
        \texttt{row-}$r$\texttt{.in}$s$, де $i = r + 4s$,
        \texttt{delay} $= 600$;
  \item створюються 16~вузлів \texttt{twiddle}
        (\texttt{tw-}$r$\texttt{-}$k_1$) з параметром
        $\{N: 16,\; k: r \cdot k_1\}$ і \texttt{latency} $= 200$;
  \item додаються 16~дуг \texttt{row-}$r$\texttt{.out}$k_1 \to$
        \texttt{tw-}$r$\texttt{-}$k_1$\texttt{.in};
  \item створюються 4~вузли \texttt{dft4} другого ярусу
        (\texttt{col-0} \ldots \texttt{col-3});
  \item додаються 16~дуг \texttt{tw-}$r$\texttt{-}$k_1$\texttt{.out}
        $\to$ \texttt{col-}$k_1$\texttt{.in}$r$;
  \item створюються 16~вузлів \texttt{sink}
        (\texttt{snk0} \ldots \texttt{snk15});
  \item додаються 16~дуг \texttt{col-}$k_1$\texttt{.out}$k_2 \to$
        \texttt{snk}$_{k_2 + 4 k_1}$\texttt{.in}.
\end{enumerate}

Повертається об'єкт \texttt{Graph} із 56~вузлами і 64~дугами. Функція
\texttt{layoutDFT4x4()} (модуль \texttt{layout4x4.ts}) обчислює
координати вузлів для React Flow: граф розподіляється на п'ять
вертикальних ярусів з кроком \texttt{stageGap} $= 280$ пікселів;
вузли в межах ярусу рівномірно розподіляються по осі~$y$ з кроком
\texttt{rowGap} $= 140$ пікселів.

\subsubsection*{Сховище стану simStore (Zustand)}

Сховище \texttt{simStore} реалізовано функцією \texttt{create()}
з Zustand і містить поля стану та методи-команди
(табл.~\ref{tab:store}).

\begin{table}[H]
  \caption{Ключові поля та методи Zustand-сховища}
  \label{tab:store}
  \centering
  \small
  \begin{tabular}{p{3.2cm}p{3.2cm}p{5.8cm}}
    \toprule
    Поле / метод & Тип & Призначення \\
    \midrule
    \texttt{graph} & \texttt{Graph} & Поточна топологія DFG \\
    \texttt{runtime} & \texttt{DataflowRuntime\textbar null} & Ядро симуляції \\
    \texttt{running} & \texttt{boolean} & Стан rAF-циклу \\
    \texttt{speed} & \texttt{number} & Множник швидкості $[0{,}1; 10]$ \\
    \texttt{simTime} & \texttt{number} & Поточний логічний час \\
    \texttt{tokensByEdge} & \texttt{Record<id, Token[]>} & Токени для анімації \\
    \texttt{sinks} & \texttt{Record<id, Complex>} & Спектральні відліки $X(k)$ \\
    \texttt{nodeActivity} & \texttt{Record<id, number>} & Теплова карта $a_i$ \\
    \texttt{metrics} & \texttt{MetricsState} & EMA-метрики \\
    \texttt{mismatches} & \texttt{Record<id, boolean>} & Результат верифікації \\
    \texttt{setGraph(g)} & \texttt{(Graph) => void} & Ініціалізація нового графа і \texttt{runtime} \\
    \texttt{injectData(v)} & \texttt{(Complex[16]) => void} & Введення вхідного вектора \\
    \texttt{start()/stop()/step()} & \texttt{() => void} & Керування rAF-циклом \\
    \texttt{applyPreset(k)} & \texttt{(PresetKind) => void} & Вбудовані пресети вхідних сигналів \\
    \bottomrule
  \end{tabular}
\end{table}

Компоненти UI підписуються на окремі зрізи стану через хук
\texttt{useSimStore(selector)} з мемоїзацією: наприклад, компонент
\texttt{Metrics} підписується лише на поле \texttt{metrics} і
перемальовується тільки при його зміні, що забезпечує стабільні
60~fps (NF3).

\subsubsection*{Анімація токенів (TokenEdge)}

Компонент \texttt{TokenEdge} (модуль \texttt{src/graph/\allowbreak TokenEdge.tsx})
відображає кожну дугу як анімовану стрілку з кружечками-токенами.
Для кожного токена у \texttt{store.tokensByEdge[edgeId]} обчислюється
відносний прогрес руху:

\begin{equation}
  \mathrm{progress} = \frac{\texttt{simTime} - \mathrm{token}.t_0}
                           {\mathrm{token}.\texttt{delay}}
                    \in [0, 1],
  \label{eq:progress}
\end{equation}

де $t_0$ --- момент виникнення токена на дузі. Поточна координата
визначається лінійною інтерполяцією між центрами вузлів-кінців:

\begin{equation}
  \mathbf{r}(\mathrm{progress}) =
    (1 - \mathrm{progress}) \cdot \mathbf{r}_{\text{src}}
    + \mathrm{progress} \cdot \mathbf{r}_{\text{tgt}}.
  \label{eq:lerp}
\end{equation}

Колір кружечка кодує тип вузла-джерела, значення $\{re, im\}$
виводиться всередині кружечка для прозорого спостереження.

\subsubsection*{Теплова карта активності вузлів}

Теплова карта реалізується у \texttt{simStore} через механізм
експоненціального затухання~(\ref{eq:heat_decay}) з підрозд.~2.6.
Функція \texttt{bumpActivity(nodeId)} викликається з обробника
події \texttt{onFire} та інкрементує лічильник відповідного вузла.
У кожному кадрі анімаційного циклу викликається
\texttt{decayActivity()}, що множить усі лічильники на
$\alpha = 0{,}92$. Поточне значення $a_i$ відображається як
яскравість заливки відповідного вузла через CSS-клас, що
обирається у модулі \texttt{\_shared.ts}.

\subsubsection*{Асинхронні обчислення у Web Worker}

Для уникнення блокування головного потоку при масштабуванні на
більші $N$ реалізовано Web Worker у файлі
\texttt{src/workers/fft.worker.ts}. Через бібліотеку Comlink
функція \texttt{applyTwiddle(tokens, factors)} приймає масив
токенів і масив поворотних множників, виконує пакетне комплексне
множення у фоновому потоці та повертає результат. У базовій
конфігурації $N = 16$ Web Worker не задіяний; він зарезервований
для подальшого розширення.

\subsubsection*{Цикл requestAnimationFrame у GraphView}

Головний анімаційний цикл запускається у хуку \texttt{useEffect}
компонента \texttt{GraphView} після ініціалізації \texttt{runtime}.
На кожному кадрі:

\begin{enumerate}
  \item якщо \texttt{store.running = true}, викликається
        \texttt{tickSimTime()}, потім
        \texttt{runtime.tick(dt, maxFires)}; у режимі
        \texttt{'single-fire'} \texttt{maxFires} $= 1$, у режимі
        \texttt{'run'} --- \texttt{Infinity};
  \item викликаються \texttt{sampleQueues()} та
        \texttt{decayActivity()} для оновлення метрик і теплової
        карти;
  \item щосекунди агрегуються лічильники \texttt{firesInWindow} для
        обчислення пропускної здатності;
  \item цикл перезапускається через \texttt{requestAnimationFrame}.
\end{enumerate}

Такий дизайн гарантує, що ядро симуляції виконує один детермінований
крок на один кадр браузера, а UI оновлюється з частотою 60~Гц.


% -------------------------------------------------
\subsection{Обробка виняткових ситуацій та забезпечення надійності}

Оскільки симулятор не взаємодіє з мережею, серверами чи файловою
системою, класичні джерела збоїв (мережеві тайм-аути, транзакційні
відмови, I/O-помилки) для нього не характерні. Водночас існують
специфічні категорії помилок, пов'язаних з валідацією вхідних даних,
цілісністю графа симуляції та обмеженнями браузерного середовища.

\subsubsection*{Класифікація типових помилок}

Помилки, що можуть виникнути у симуляторі, зведено у
табл.~\ref{tab:errors}.

\begin{table}[H]
  \caption{Класифікація помилок симулятора та підходи до їх обробки}
  \label{tab:errors}
  \centering
  \small
  \begin{tabular}{p{2.8cm}p{4.0cm}p{5.6cm}}
    \toprule
    Категорія & Приклади & Обробка \\
    \midrule
    Валідаційні & Введення не-числового значення у \texttt{SourceNode}, множник швидкості поза $[0{,}1; 10]$ & Clamp до допустимого діапазону, підсвічування червоним \\
    Транзієнтні & Пропуск кадру при перевантаженні CPU & \texttt{maxFires} обмежує кількість активацій за один \texttt{tick()} \\
    Структурні & Відсутність вхідного токена на порті вузла & Вузол не активується (правило готовності); статичний аналіз графа виключає deadlock \\
    Чисельні & Втрата точності при $|X(k)| > 10^{12}$ & Використання IEEE 754 double; попередження при виході за рамки \\
    Браузерні & Відсутність Web Worker API & Graceful degradation: обчислення у головному потоці \\
    \bottomrule
  \end{tabular}
\end{table}

\subsubsection*{Валідація вхідних даних}

Компонент \texttt{SourceNode} використовує повзунки з фіксованими
межами $[-10; 10]$ для \texttt{re} та \texttt{im}, що виключає
введення нечислових значень. Компонент \texttt{Controls} аналогічно
обмежує множник \texttt{speed} діапазоном $[0{,}1; 10]$ за допомогою
\texttt{Math.max(0.1, Math.min(10, value))}. Параметри затримок
\texttt{latency} і \texttt{delay} валідуються на додатність і
цілочисельність; у разі порушення викидається виняток
\texttt{InvalidGraphError} з описовою причиною.

\subsubsection*{Відсутність ризику deadlock у DFG}

Оскільки граф симулятора є статичним орієнтованим ациклічним
графом (DAG) з детермінованою SDF-семантикою (підрозд.~1.2), ризик
взаємного блокування принципово відсутній: відсутність циклів
гарантує, що готовність вузлів поступово поширюється від
\texttt{source} до \texttt{sink}. Ця властивість перевіряється на
етапі побудови графа функцією \texttt{validateDAG(graph)} через
топологічне сортування методом Кана.

\subsubsection*{Логування та трасування}

Усі значущі події ядра (\texttt{onToken}, \texttt{onFire},
\texttt{onOutput}) акумулюються у \texttt{store.eventLog} із
часовими мітками \texttt{simTime} та \texttt{performance.now()}.
Компонент \texttt{TerminalOutput} відображає журнал у реальному
часі з можливістю фільтрації за типом події. Експорт журналу у
файл (CSV/JSON) реалізований через Web API
\texttt{Blob + URL.createObjectURL} без передачі на сервер.

Рівні логування (\texttt{debug}, \texttt{info}, \texttt{warn},
\texttt{error}) налаштовуються через прапорець
\texttt{store.logLevel}; у production-збірці рівень
\texttt{debug} вимкнений за замовчуванням.

\subsubsection*{Graceful degradation та оптимізація}

\textbf{Fallback при відсутності Web Worker.} Якщо браузер не
підтримує Web Workers, функція \texttt{applyTwiddle()} виконується
синхронно у головному потоці; для $N = 16$ це не впливає на
частоту кадрів.

\textbf{Оптимізація перерендерингу.} Селектори Zustand
мемоїзовані функцією \texttt{shallow}, що запобігає зайвим
перемалюванням компонентів при зміні несуттєвих полів сховища.
Важкі розрахунки (наприклад, обчислення координат токенів на
ребрах) виконуються лише для видимої області через перевірку
\texttt{ReactFlow.viewportBounds}.

\textbf{Профілювання.} Для розробника у режимі debug доступна панель
з часовими показниками: тривалість \texttt{tick()}, тривалість
рендеру React, інтервали між кадрами. Середня тривалість одного
кроку \texttt{tick()} для $N = 16$ становить $< 1$~мс, що залишає
понад 15~мс запасу у бюджеті кадру (16.67~мс при 60~fps).


% -------------------------------------------------
\subsection{Розгортання системи}

Розгортання симулятора спрощене порівняно з класичними
веб-застосунками завдяки відсутності серверної частини та бази
даних. Зібраний статичний бандл може бути розміщений на будь-якому
файловому хостингу, що віддає статичні ресурси.

\subsubsection*{Середовище розгортання}

Цільовим середовищем є \textbf{статичний веб-хостинг}: GitHub Pages,
Vercel (статичний режим), Netlify або локальний HTTP-сервер
розробника. Серверна інфраструктура, бази даних, сесійне сховище
та автентифікація не застосовуються.

\subsubsection*{Процедура збірки}

Збірка виконується стандартною командою Next.js:

\begin{figure}[H]
\begin{verbatim}
$ pnpm install            # встановлення залежностей
$ pnpm run build          # next build -> .next/
$ pnpm run export         # next export -> out/
$ pnpm run lint           # ESLint
$ pnpm test               # Jest-тести
\end{verbatim}
\caption{Команди збірки та перевірки проєкту. Джерело: розроблено
         автором.}
\label{lst:build}
\end{figure}

У результаті у директорії \texttt{out/} формується статичний
бандл: HTML-сторінка, JavaScript-бандл із коду застосунку (через
Turbopack/webpack із оптимізаціями tree-shaking і minification) та
CSS-файли. Розмір gzip-бандла становить орієнтовно 250~КБ.

\subsubsection*{CI/CD-пайплайн (опційно)}

Для автоматизації перевірок на кожному push до основної гілки
налаштовано GitHub Actions із workflow \texttt{.github/workflows/ci.yml}:

\begin{itemize}
  \item job \texttt{lint} --- запуск ESLint і TypeScript-компіляції
        у режимі \texttt{--noEmit} (перевірка типів без генерації
        файлів);
  \item job \texttt{test} --- запуск Jest-тестів з генерацією звіту
        про покриття;
  \item job \texttt{build} --- перевірка, що збірка проходить без
        помилок.
\end{itemize}

Стратегія релізу --- \textit{rolling}: нова статична збірка замінює
попередню при успішному проходженні CI. Оскільки застосунок не має
сесійного стану, жодних процедур міграції бази даних чи інвалідації
кешу не потрібно.

\subsubsection*{Резервне копіювання та відтворюваність}

Оскільки вихідний код є єдиним артефактом, що підлягає збереженню,
резервне копіювання реалізоване через Git: повна історія зберігається
у віддаленому репозиторії GitHub. Детермінованість збірки гарантується
файлом \texttt{pnpm-lock.yaml}, що фіксує точні версії всіх
залежностей. Будь-який виконавець може відтворити продуктову збірку
командою \texttt{pnpm install \&\& pnpm run build} на машині з
Node.js $\ge 18$.


% -------------------------------------------------
\subsection{Тестування програмного забезпечення}

Тестування на етапі розробки охоплює модульні тести найкритичніших
компонентів і базові інтеграційні сценарії. Комплексне тестування
готової системи, порівняння з еталонними результатами та аналіз
продуктивності розглядаються у розділі~4.

\subsubsection*{Фреймворк тестування}

Для модульного тестування використовується \textbf{Jest~29.x} ---
загальновживаний фреймворк JavaScript/TypeScript-тестів з
вбудованим runner'ом, матчерами та покриттям. Конфігурація
винесена у файл \texttt{jest.config.ts}, де налаштовано:

\begin{itemize}
  \item середовище --- \texttt{node} (для модулів \texttt{core/} та
        \texttt{utils/}) та \texttt{jsdom} (для UI-тестів);
  \item трансформатор --- \texttt{ts-jest} з \texttt{tsconfig.test.json};
  \item поріг покриття --- не менше 70~\% для тверджень,
        віток і функцій у директоріях \texttt{src/core/} та
        \texttt{src/utils/}.
\end{itemize}

\subsubsection*{Перелік тест-сьютів}

Реалізовані тест-сьюти та цільові модулі наведено у
табл.~\ref{tab:tests}.

\begin{table}[H]
  \caption{Модульні тести симулятора}
  \label{tab:tests}
  \centering
  \small
  \begin{tabular}{p{3.2cm}p{4.0cm}p{5.8cm}}
    \toprule
    Тест-файл & Цільовий модуль & Сценарії \\
    \midrule
    \texttt{math.test.ts} & \texttt{utils/math.ts} & Коректність \texttt{cplxAdd}, \texttt{cplxSub}, \texttt{cplxMulJ}; правильність \texttt{computeDFT4} на референтних входах \\
    \texttt{buffers.test.ts} & \texttt{core/buffers.ts} & FIFO-порядок, поведінка \texttt{peek} на пустій черзі, розмір \\
    \texttt{scheduler.test.ts} & \texttt{core/scheduler.ts} & Правило активації, коректність \texttt{tick()}, дотримання \texttt{maxFires}, детермінованість \\
    \texttt{generate.test.ts} & \texttt{graph/generateDFT4x4.ts} & Кількість вузлів (56) і дуг (64), коректність з'єднань \\
    \texttt{compare.test.ts} & \texttt{utils/compare.ts} & Збіг симулятора з fft.js на імпульсі та синусоїді з $\varepsilon \le 10^{-9}$ \\
    \texttt{store.test.ts} & \texttt{store/simStore.ts} & Скидання стану, \texttt{injectData}, \texttt{applyPreset} \\
    \bottomrule
  \end{tabular}
\end{table}

\subsubsection*{Приклад тест-кейса}

Лістинг~\ref{lst:test_example} наводить приклад тесту коректності
функції \texttt{computeDFT4} для імпульсного входу $(1, 0, 0, 0)$:
очікуваний результат --- $(1, 1, 1, 1)$, оскільки всі поворотні
множники в сумі дають одиничні коефіцієнти.

\begin{figure}[H]
\begin{verbatim}
describe('computeDFT4', () => {
  it('returns (1,1,1,1) for unit impulse (1,0,0,0)', () => {
    const input: Complex[] = [
      { re: 1, im: 0 }, { re: 0, im: 0 },
      { re: 0, im: 0 }, { re: 0, im: 0 }
    ];
    const out = computeDFT4(input);
    for (const z of out) {
      expect(z.re).toBeCloseTo(1, 12);
      expect(z.im).toBeCloseTo(0, 12);
    }
  });
});
\end{verbatim}
\caption{Приклад модульного тесту функції \texttt{computeDFT4}.
         Джерело: розроблено автором.}
\label{lst:test_example}
\end{figure}

\subsubsection*{Рівень покриття коду}

Поточний рівень покриття, отриманий командою \texttt{pnpm test
--coverage}, становить:

\begin{itemize}
  \item \texttt{src/core/} --- 85~\% statements, 78~\% branches,
        92~\% functions;
  \item \texttt{src/utils/} --- 91~\% statements, 84~\% branches,
        100~\% functions;
  \item \texttt{src/graph/generateDFT4x4.ts} --- 100~\% statements.
\end{itemize}

Значення перевищують цільовий поріг 70~\%, визначений у
\texttt{jest.config.ts}. UI-компоненти (\texttt{src/components/})
охоплені лише дим-тестами на монтування та не враховуються у
обчисленні покриття, оскільки їх коректність підтверджується
ручним тестуванням у розділі~4.

\subsubsection*{Автоматизоване виконання тестів у CI (опційно)}

Тести автоматично виконуються у GitHub Actions на кожному push до
\texttt{main} або pull request. Конфігурація workflow
\texttt{.github/workflows/ci.yml} містить кроки установлення
залежностей, компіляції TypeScript, запуску ESLint і Jest. Звіт
про покриття у форматі HTML публікується як артефакт збірки; у
README додано badge з поточним відсотком покриття.


% -------------------------------------------------
\subsection*{Висновки до розділу~3}

У третьому розділі подано програмну реалізацію симулятора, що
відповідає архітектурі, спроєктованій у розділі~2.

\emph{Описано технологічний стек:} TypeScript~5 як статично
типізована мова реалізації, Next.js~14 як фреймворк статичної SPA,
React~18 і React Flow~12 для реактивного UI та графової
візуалізації, Zustand~4 для централізованого сховища стану,
Jest~29 для модульних тестів, ESLint/Prettier для забезпечення
стилю. Усі обрані бібліотеки мають відкриті ліцензії MIT або
Apache~2.0, що задовольняє вимогу NF6.

\emph{Організовано вихідний код} за принципом розподілу за шарами
відповідальності: ядро симуляції (\texttt{src/core/}) без
залежностей від UI, генерація та візуалізація графа
(\texttt{src/graph/}), UI-компоненти (\texttt{src/components/}),
сховище стану (\texttt{src/store/}), утиліти (\texttt{src/utils/})
та Web Workers (\texttt{src/workers/}). Повну відповідність
архітектурних компонентів з підрозд.~2.3 до модулів коду подано
у табл.~\ref{tab:arch_to_code}.

\emph{Реалізовано ключові алгоритми:} клас \texttt{DataflowRuntime}
з методом \texttt{tick()} у чотири кроки
(доставка~(\ref{eq:delivery}) $\to$ перевірка готовності $\to$
активація $\to$ розміщення вихідних токенів~(\ref{eq:emit_time}));
функція \texttt{generateDFT4x4()} як Factory для побудови DFG із
56~вузлів і 64~дуг; сховище \texttt{simStore} на Zustand із
підпискою UI на окремі зрізи стану; компоненти \texttt{TokenEdge}
та теплової карти з інтерполяцією~(\ref{eq:progress})--(\ref{eq:lerp})
і експоненціальним затуханням активності.

\emph{Розглянуто обробку виняткових ситуацій:} класифікацію помилок
за категоріями (валідаційні, транзієнтні, структурні, чисельні,
браузерні), підхід до валідації вхідних даних через обмеження UI,
доведено відсутність ризику deadlock у DAG-структурі симулятора,
описано механізми логування та graceful degradation при
відсутності Web Worker API.

\emph{Визначено процедуру розгортання} як статичної SPA на довільному
файловому хостингу з підтримкою HTTPS. Налаштовано (опційно)
GitHub Actions для автоматичної перевірки типів, лінтера та тестів
на кожному push.

\emph{Проведено модульне тестування} у Jest з покриттям 85~\% для
\texttt{src/core/} та 91~\% для \texttt{src/utils/}, що перевищує
цільовий поріг 70~\%. Реалізовано шість тест-сьютів, що охоплюють
комплексну арифметику, FIFO-черги, планувальник, генератор графа,
верифікацію та сховище стану. Комплексна експериментальна оцінка
якості роботи системи --- точність відносно еталонного fft.js,
продуктивність, аналіз критичного шляху --- виноситься до розділу~4.